\chapter{Interfaz interactiva y simulación (frontend en Vue.js)}

\section{Arquitectura del frontend}

La interfaz carga Vue~3 desde CDN y monta \texttt{\#app} con la API de opciones. Estado reactivo principal:

\begin{table}[H]
  \centering
  \caption{Estado reactivo de la simulación Vue}
  \begin{tabular}{@{}lp{9cm}@{}}
    \toprule
    \textbf{Propiedad} & \textbf{Función} \\
    \midrule
    \texttt{root} & Snapshot del árbol desde API \\
    \texttt{stats} & Contadores: nodos, altura, $q$, rebuilds \\
    \texttt{activeNodes} & Mapa id $\rightarrow$ tipo de animación CSS \\
    \texttt{stepText} & Mensaje pedagógico del paso actual \\
    \bottomrule
  \end{tabular}
\end{table}

\section{Flujo de animación del scapegoat}

Tras \texttt{POST /api/products}, la función \texttt{playInsertTrace}:
\begin{enumerate}
  \item Recorre \texttt{trace.path} nodo a nodo (\texttt{playPath}).
  \item Muestra profundidad vs.\ \texttt{maxDepth}.
  \item Si \texttt{trace.scapegoat} existe: resalta el nodo culpable y luego \texttt{rebuiltKeys}.
\end{enumerate}

\begin{figure}[H]
  \centering
  \begin{tikzpicture}[
    state/.style={rectangle, rounded corners, draw, minimum width=2.4cm, minimum height=0.7cm, align=center, font=\footnotesize},
    arrow/.style={-{Stealth[length=2mm]}, thick}
  ]
    \node[state] (a) {Cargar árbol};
    \node[state, below=0.7cm of a] (b) {Insertar};
    \node[state, below=0.7cm of b] (c) {Animar camino};
    \node[state, below=0.7cm of c] (d) {¿Scapegoat?};
    \node[state, below left=0.8cm and 1cm of d] (e) {Fin};
    \node[state, below right=0.8cm and 1cm of d] (f) {Rebuild visual};
    \draw[arrow] (a) -- (b);
    \draw[arrow] (b) -- (c);
    \draw[arrow] (c) -- (d);
    \draw[arrow] (d) -- node[left, font=\tiny]{No} (e);
    \draw[arrow] (d) -- node[right, font=\tiny]{Sí} (f);
    \draw[arrow] (f) |- (a);
  \end{tikzpicture}
  \caption{Ciclo de animación en Vue.js}
  \label{fig:vue-anim}
\end{figure}

\section{Uso de inteligencia artificial}

El frontend recibió asistencia de herramientas de IA generativa para layout HTML/CSS y animaciones. El equipo revisó manualmente \texttt{findPath}, \texttt{renderTree} y la integración con \texttt{InsertTrace}. El algoritmo del \Scapegoat{} y la capa de persistencia fueron implementados sin asistencia automatizada en su lógica central.
